\begin{apendicesenv}

\apendiceitem{Especificação XML do Diagrama de Componentes no UMLet}{apendice:uml}

O trecho XML reproduzido na Listagem \ref{lst:umlet} corresponde ao arquivo \texttt{diagrama\_de\_componentes.uxf}, produzido pelos autores na ferramenta UMLet para modelar os componentes \textit{Usuários}, \textit{Fila}, \textit{Pedido de Música} e \textit{Reprodutor}.

\begin{lstlisting}[language=XML, caption={Estrutura do arquivo diagrama\_de\_componentes.uxf (UMLet).}, label={lst:umlet}]
<?xml version="1.0" encoding="UTF-8" standalone="no"?>
<diagram program="umlet" version="15.1">
  <zoom_level>10</zoom_level>
  <element>
    <id>UMLGeneric</id>
    <coordinates><x>250</x><y>80</y><w>430</w><h>110</h></coordinates>
    <panel_attributes>symbol=component
Usuarios</panel_attributes>
  </element>
  <element>
    <id>UMLGeneric</id>
    <coordinates><x>250</x><y>280</y><w>430</w><h>100</h></coordinates>
    <panel_attributes>symbol=component
Fila</panel_attributes>
  </element>
  <element>
    <id>UMLGeneric</id>
    <coordinates><x>160</x><y>530</y><w>200</w><h>100</h></coordinates>
    <panel_attributes>symbol=component
Pedido de Musica</panel_attributes>
  </element>
  <element>
    <id>UMLGeneric</id>
    <coordinates><x>520</x><y>520</y><w>200</w><h>100</h></coordinates>
    <panel_attributes>symbol=component
Reprodutor</panel_attributes>
  </element>
</diagram>
\end{lstlisting}

\apendiceitem{Trechos de Código-Fonte do MVP (Electron e Validação da Fila)}{apendice:codigo}

A Listagem \ref{lst:electron_main} apresenta o script principal de inicialização do \textbf{Karaokê TV Player} em Electron (\texttt{main.js}), incluindo a configuração da janela dedicada de exibição e o isolamento de contexto.

\begin{lstlisting}[language=Java, caption={Inicialização do reprodutor desktop Karaokê TV Player em Electron (main.js).}, label={lst:electron_main}]
const { app, BrowserWindow } = require('electron');
const path = require('path');

function createPlayerWindow() {
  const playerWindow = new BrowserWindow({
    width: 1280,
    height: 720,
    title: 'Karaoke TV Player - Apache Gru Rock Bar',
    autoHideMenuBar: true,
    webPreferences: {
      preload: path.join(__dirname, 'preload.js'),
      contextIsolation: true,
      nodeIntegration: false,
      sandbox: true
    }
  });

  const playerUrl = process.env.ELEC_PLAYER_URL || 'http://localhost:5173/player';
  playerWindow.loadURL(playerUrl);
}

app.whenReady().then(() => {
  createPlayerWindow();
});
\end{lstlisting}

A Listagem \ref{lst:youtube_parser} demonstra a função de extração de identificadores de vídeo do YouTube e validação das regras RN00 e RN01 antes da inserção na fila.

\begin{lstlisting}[language=Java, caption={Extração de ID do YouTube e validação anti-duplicidade na fila (RN00 e RN01).}, label={lst:youtube_parser}]
export function extractYouTubeId(url) {
  if (!url || typeof url !== 'string') return null;
  const regex = /(?:youtube\.com\/(?:[^\/]+\/.+\/|(?:v|e(?:mbed)?)\/|.*[?&]v=)|youtu\.be\/)([^"&?\/\s]{11})/i;
  const match = url.trim().match(regex);
  return match ? match[1] : null;
}

export function validateQueueInsertion(videoUrl, customTitle, activeQueue, currentSong) {
  const videoId = extractYouTubeId(videoUrl);
  if (!videoId) {
    return { valid: false, error: 'RN00: Link do YouTube invalido ou nao suportado.' };
  }
  const alreadyPlaying = currentSong && currentSong.videoId === videoId;
  const alreadyInQueue = activeQueue.some((item) => item.videoId === videoId);
  if (alreadyPlaying || alreadyInQueue) {
    return { valid: false, error: 'RN01: Esta musica ja esta na fila ou em reproducao.' };
  }
  return {
    valid: true,
    songItem: {
      videoId,
      title: customTitle?.trim() || `Musica (${videoId})`,
      status: 'na_fila',
      requestedAt: new Date().toISOString()
    }
  };
}
\end{lstlisting}

\end{apendicesenv}

%% -------------------------------------------------------------------------- %%
